% ──────────────────────────────────────────────────────────────────────────────
% Presentación: Detección y Eliminación de Sellos — Resultados y Contraste
% Autor: Manuel Cruz Garrote | Universidad del Rosario | Mayo 2026
% ──────────────────────────────────────────────────────────────────────────────
\documentclass[aspectratio=169, 11pt]{beamer}

% ─── PAQUETES ─────────────────────────────────────────────────────────────────
\usepackage[T1]{fontenc}
\usepackage[utf8]{inputenc}
\usepackage[spanish]{babel}
\usepackage{lmodern}
\usepackage{booktabs}
\usepackage{graphicx}
\usepackage{xcolor}
\usepackage{tikz}
\usepackage{enumitem}
\usepackage{array}
\usepackage{colortbl}
\usepackage{multirow}
\usepackage{pgfplots}
\pgfplotsset{compat=1.18}
\usetikzlibrary{arrows.meta, shapes.rounded, positioning, fit, backgrounds}

% ─── TEMA ─────────────────────────────────────────────────────────────────────
\usetheme{metropolis}
\setbeamertemplate{navigation symbols}{}

% ─── PALETA DE COLORES ────────────────────────────────────────────────────────
\definecolor{PriRed}{RGB}{192,0,0}
\definecolor{AccBlue}{RGB}{31,73,125}
\definecolor{AccGreen}{RGB}{0,112,60}
\definecolor{AccAmber}{RGB}{180,120,0}
\definecolor{AccTeal}{RGB}{0,110,120}
\definecolor{MidGray}{RGB}{100,100,100}
\definecolor{LightBg}{RGB}{248,248,248}
\definecolor{DarkBg}{RGB}{20,20,20}
\definecolor{OldRed}{RGB}{160,30,30}
\definecolor{NewGreen}{RGB}{0,120,70}
\definecolor{RowA}{RGB}{245,245,255}
\definecolor{RowB}{RGB}{240,255,240}

\setbeamercolor{frametitle}{bg=AccBlue, fg=white}
\setbeamercolor{progress bar}{fg=PriRed}
\setbeamercolor{alerted text}{fg=PriRed}
\setbeamercolor{example text}{fg=AccGreen}

% ─── COMANDOS UTILITARIOS ─────────────────────────────────────────────────────
\newcommand{\rojo}[1]{\textcolor{PriRed}{#1}}
\newcommand{\azul}[1]{\textcolor{AccBlue}{#1}}
\newcommand{\verde}[1]{\textcolor{AccGreen}{#1}}
\newcommand{\amber}[1]{\textcolor{AccAmber}{#1}}
\newcommand{\lib}[1]{\texttt{\small #1}}
\newcommand{\badge}[2]{\tikz[baseline=-0.5ex]\node[fill=#1,text=white,
  rounded corners=2pt,inner sep=2pt,font=\tiny\bfseries]{#2};}
\newcommand{\mejora}[1]{\verde{\textbf{#1}}}
\newcommand{\original}[1]{\textcolor{OldRed}{\textbf{#1}}}

% ─── TIPOGRAFÍA ───────────────────────────────────────────────────────────────
\setbeamerfont{frametitle}{size=\large, series=\bfseries}
\setbeamerfont{title}{size=\LARGE, series=\bfseries}

% ─── BLOQUES PERSONALIZADOS ───────────────────────────────────────────────────
\newenvironment{redblock}[1]{%
  \setbeamercolor{block title}{fg=white, bg=OldRed}%
  \setbeamercolor{block body}{fg=DarkBg, bg=OldRed!8}%
  \begin{block}{#1}}{\end{block}}

\newenvironment{greenblock}[1]{%
  \setbeamercolor{block title}{fg=white, bg=AccGreen}%
  \setbeamercolor{block body}{fg=DarkBg, bg=AccGreen!8}%
  \begin{block}{#1}}{\end{block}}

\newenvironment{blueblock}[1]{%
  \setbeamercolor{block title}{fg=white, bg=AccBlue}%
  \setbeamercolor{block body}{fg=DarkBg, bg=AccBlue!8}%
  \begin{block}{#1}}{\end{block}}

% ─── METADATOS ────────────────────────────────────────────────────────────────
\title{Detección y Eliminación Inteligente de Sellos\\para Mejora del OCR}
\subtitle{Contraste de Metodologías y Resultados Experimentales}
\author{Manuel Cruz Garrote}
\institute{Universidad del Rosario}
\date{Mayo 2026}

% ══════════════════════════════════════════════════════════════════════════════
\begin{document}
% ══════════════════════════════════════════════════════════════════════════════

% ─── PORTADA ──────────────────────────────────────────────────────────────────
\begin{frame}[plain]
  \maketitle
\end{frame}

% ─── ÍNDICE ───────────────────────────────────────────────────────────────────
\begin{frame}{Contenido}
  \setbeamertemplate{section in toc}[sections numbered]
  \tableofcontents[hideallsubsections]
\end{frame}

% ══════════════════════════════════════════════════════════════════════════════
\section{Contexto y Problema}
% ══════════════════════════════════════════════════════════════════════════════

\begin{frame}{El Problema: Sellos sobre Texto}
  \begin{columns}[T]
    \begin{column}{0.55\textwidth}
      \begin{itemize}[leftmargin=1.2em, itemsep=8pt]
        \item Los documentos oficiales en español presentan
              \rojo{\textbf{sellos translúcidos}} que se superponen
              sobre el texto base.
        \item Las tasas de error OCR en la región cubierta son
              \rojo{\textbf{hasta 3$\times$ superiores}} al resto del documento.
        \item El problema afecta notarías, registros civiles y
              expedientes judiciales digitalizados.
      \end{itemize}
      \vspace{8pt}
      \begin{exampleblock}{Objetivo central}
        Reducir el CER en \textbf{al menos 5\,pp} respecto al
        baseline sin preprocesamiento.
      \end{exampleblock}
    \end{column}
    \begin{column}{0.42\textwidth}
      \centering
      \begin{tikzpicture}
        % Documento
        \fill[white, draw=MidGray, rounded corners=4pt] (0,0) rectangle (3.8,4.8);
        % Líneas de texto simuladas
        \foreach \y in {0.4,0.8,1.2,1.6,2.0,2.4,2.8,3.2,3.6,4.0,4.4}
          \fill[MidGray!40, rounded corners=1pt] (0.2,\y-0.05) rectangle (3.6,\y+0.1);
        % Sello translúcido
        \fill[red!30, opacity=0.6] (1.0,1.5) circle (1.1);
        \draw[red!60, line width=1.2pt] (1.0,1.5) circle (1.1);
        \draw[red!60, line width=0.6pt] (1.0,1.5) circle (0.85);
        \node[red!70, font=\tiny\bfseries, rotate=15] at (1.0,1.5) {SELLO};
        % Etiqueta
        \node[font=\tiny, text=MidGray, below] at (1.9,0) {Documento con sello translúcido};
      \end{tikzpicture}
    \end{column}
  \end{columns}
\end{frame}

% ══════════════════════════════════════════════════════════════════════════════
\section{Contraste de Metodologías}
% ══════════════════════════════════════════════════════════════════════════════

\begin{frame}{Evolución: De la Propuesta Inicial a la Implementación}
  \begin{columns}[T]
    \begin{column}{0.48\textwidth}
      \begin{redblock}{Propuesta Original (LaTeX)}
        \scriptsize
        \begin{itemize}[leftmargin=1em, itemsep=4pt]
          \item \textbf{Detección:} YOLOv8 entrenado en dataset COCO
          \item \textbf{Segmentación:} Umbralización HSV fija
          \item \textbf{Inpainting:} LaMa (large mask, GPU requerida)
          \item \textbf{OCR:} PaddleOCR / Tesseract v5
          \item \textbf{Eval:} \lib{jiwer.cer()} + IoU + SSIM/FID
          \item \textbf{Corrección:} pyspellchecker básico
          \item \textbf{Datos:} 100 docs benchmark humana
        \end{itemize}
      \end{redblock}
    \end{column}
    \begin{column}{0.48\textwidth}
      \begin{greenblock}{Implementación Actual (Pipeline)}
        \scriptsize
        \begin{itemize}[leftmargin=1em, itemsep=4pt]
          \item \textbf{Detección:} Clasificación cromática (5 clases RGB)
          \item \textbf{Segmentación:} Horizontal + block sweep filters
          \item \textbf{Inpainting:} Híbrido TELEA + Navier-Stokes
          \item \textbf{OCR:} EasyOCR + corrección Levenshtein
          \item \textbf{Eval:} LevenshteinCalculator propio (CER/WER)
          \item \textbf{Corrección:} spaCy NER + textdistance
          \item \textbf{Datos:} 200 imágenes sintéticas + docs reales
        \end{itemize}
      \end{greenblock}
    \end{column}
  \end{columns}
\end{frame}

\begin{frame}{Tabla de Contraste Técnico: Decisiones de Diseño}
  \centering
  \footnotesize
  \begin{tabular}{@{}p{2.6cm}p{4.2cm}p{4.2cm}@{}}
    \toprule
    \textbf{Componente}
      & \original{Propuesta Original}
      & \mejora{Implementado} \\
    \midrule
    \rowcolor{RowA}
    Clasificación píxel
      & Umbralización HSV con rangos fijos
      & 5 categorías cromáticas (TEXTO, SELLO\_SAT, MIXTO, FONDO, SELLO\_TRANS) \\[4pt]
    \rowcolor{RowB}
    Filtrado espacial
      & Morfología binaria estándar
      & H-sweep + block-sweep para coherencia contextual \\[4pt]
    \rowcolor{RowA}
    Inpainting
      & LaMa (requiere GPU, >4\,GB VRAM)
      & Híbrido TELEA/NS adaptativo, solo CPU \\[4pt]
    \rowcolor{RowB}
    Motor OCR
      & PaddleOCR (es\_PP-OCRv4)
      & EasyOCR con umbral de confianza configurable \\[4pt]
    \rowcolor{RowA}
    Post-corrección
      & pyspellchecker (diccionario estático)
      & Levenshtein $\leq 2$ + NER spaCy (protege entidades) \\[4pt]
    \rowcolor{RowB}
    Evaluación CER/WER
      & jiwer externo (caja negra)
      & LevenshteinCalculator propio + normalización de texto \\[4pt]
    \rowcolor{RowA}
    Portabilidad
      & GPU requerida para LaMa
      & 100\,\% CPU — ejecutable en Colab gratuito \\
    \bottomrule
  \end{tabular}
\end{frame}

\begin{frame}{Por Qué Se Abandonó la Ruta YOLO}
  \begin{columns}[T]
    \begin{column}{0.48\textwidth}
      \begin{redblock}{Limitaciones de YOLOv8}
        \begin{itemize}[leftmargin=1em, itemsep=5pt]
          \item Requiere \textbf{anotaciones manuales} (bounding boxes)
                para cada sello — costoso
          \item Falla en \textbf{sellos translúcidos} sin contraste nítido
          \item Necesita GPU para inferencia en tiempo real
          \item Modelo entrenado: BoxF1 = 0.62 en validación\\
                (bajo para producción)
          \item No generaliza a sellos de color arbitrario
        \end{itemize}
      \end{redblock}
    \end{column}
    \begin{column}{0.48\textwidth}
      \begin{greenblock}{Ventajas del Enfoque Cromático}
        \begin{itemize}[leftmargin=1em, itemsep=5pt]
          \item \textbf{Sin anotaciones}: basado en propiedades físicas del color
          \item Detecta sellos translúcidos por saturación diferencial
          \item Funciona en CPU — latencia $\sim$8\,s por imagen
          \item Configurable: solo 3 umbrales (sat, val, ratio)
          \item Generaliza por \textbf{categoría de píxel}, no por forma
        \end{itemize}
      \end{greenblock}
    \end{column}
  \end{columns}
\end{frame}

% ══════════════════════════════════════════════════════════════════════════════
\section{Arquitectura del Pipeline Actual}
% ══════════════════════════════════════════════════════════════════════════════

\begin{frame}{Flujo del Pipeline — 3 Fases Principales}
  \centering
  \begin{tikzpicture}[
    node distance=0.4cm and 0.35cm,
    box/.style={rectangle, rounded corners=5pt, minimum width=2.0cm,
                minimum height=1.0cm, align=center, font=\scriptsize\bfseries,
                text=white, inner sep=5pt},
    subbox/.style={rectangle, rounded corners=3pt, minimum width=1.8cm,
                   minimum height=0.6cm, align=center, font=\tiny, inner sep=3pt},
    arr/.style={-{Stealth[length=5pt]}, thick}
  ]

  % Nodos principales
  \node[box, fill=AccBlue] (input) {Documento\\Escaneado};
  \node[box, fill=OldRed, right=0.7cm of input] (p3a) {Fase 3A\\Stamp Removal};
  \node[box, fill=AccAmber, right=0.7cm of p3a] (p4) {Fase 4\\OCR + NLP};
  \node[box, fill=AccGreen, right=0.7cm of p4] (p5) {Fase 5\\Evaluación};
  \node[box, fill=MidGray, right=0.7cm of p5] (out) {Resultados\\JSON + TXT};

  % Flechas principales
  \draw[arr] (input) -- (p3a);
  \draw[arr] (p3a) -- (p4);
  \draw[arr] (p4) -- (p5);
  \draw[arr] (p5) -- (out);

  % Subnodos Fase 3A
  \node[subbox, fill=OldRed!20, below=0.3cm of p3a, xshift=-0.5cm] (c1)
    {\textcolor{OldRed}{Clasif.\ cromática}\\5 categorías};
  \node[subbox, fill=OldRed!20, below=0.3cm of p3a, xshift=0.9cm] (c2)
    {\textcolor{OldRed}{H/Block Sweep}\\+ Inpainting};

  % Subnodos Fase 4
  \node[subbox, fill=AccAmber!20, below=0.3cm of p4, xshift=-0.5cm] (c3)
    {\textcolor{AccAmber}{EasyOCR}\\confianza};
  \node[subbox, fill=AccAmber!20, below=0.3cm of p4, xshift=0.9cm] (c4)
    {\textcolor{AccAmber}{Levenshtein}\\+ spaCy NER};

  % Subnodos Fase 5
  \node[subbox, fill=AccGreen!20, below=0.3cm of p5] (c5)
    {\textcolor{AccGreen}{CER / WER}\\vs.\ ground truth};

  \end{tikzpicture}

  \vspace{10pt}
  \begin{columns}[T]
    \begin{column}{0.32\textwidth}
      \scriptsize\centering
      \badge{OldRed}{Fase 3A}\ \textbf{Detección + Limpieza}\\[3pt]
      Clasifica píxeles por color,\\ filtra con sweep bidireccional,\\ inpaint adaptativo
    \end{column}
    \begin{column}{0.32\textwidth}
      \scriptsize\centering
      \badge{AccAmber}{Fase 4}\ \textbf{OCR + Post-proc.}\\[3pt]
      EasyOCR español, corrección\\ ortográfica con umbral dist.\ $\leq 2$,\\ NER protege entidades
    \end{column}
    \begin{column}{0.32\textwidth}
      \scriptsize\centering
      \badge{AccGreen}{Fase 5}\ \textbf{Métricas}\\[3pt]
      CER y WER por distancia\\ Levenshtein, comparado\\ con ground truth manual
    \end{column}
  \end{columns}
\end{frame}

% ══════════════════════════════════════════════════════════════════════════════
\section{Resultados Experimentales}
% ══════════════════════════════════════════════════════════════════════════════

\begin{frame}{Resultados de Evaluación CER/WER — Documento Real}
  \begin{columns}[T]
    \begin{column}{0.48\textwidth}
      \begin{blueblock}{Condición evaluada}
        \scriptsize
        \begin{itemize}[leftmargin=1em, itemsep=4pt]
          \item \textbf{Documento:} train-00000-of-00024.parquet (muestra 84)
          \item \textbf{Ground truth:} 3{,}715 caracteres, 525 palabras
          \item \textbf{OCR extraído:} 3{,}624 caracteres, 534 palabras
          \item \textbf{Pipeline:} Fase 3A (inpainting híbrido) + Fase 4 (EasyOCR)
          \item \textbf{Corrección:} Levenshtein $\leq 2$ + spaCy NER activos
        \end{itemize}
      \end{blueblock}
    \end{column}
    \begin{column}{0.48\textwidth}
      \centering
      \vspace{6pt}

      % Métrica CER
      \begin{tikzpicture}
        \node[draw=AccGreen, line width=1.5pt, rounded corners=6pt,
              fill=AccGreen!8, minimum width=5.0cm, minimum height=1.2cm]
          (cer) {\Large\bfseries\verde{10.63\,\%}\ \ CER};
        \node[font=\scriptsize, below=2pt of cer, text=AccGreen]
          {Character Error Rate — \textbf{Excelente}};
      \end{tikzpicture}

      \vspace{8pt}

      % Métrica WER
      \begin{tikzpicture}
        \node[draw=AccAmber, line width=1.5pt, rounded corners=6pt,
              fill=AccAmber!8, minimum width=5.0cm, minimum height=1.2cm]
          (wer) {\Large\bfseries\amber{28.76\,\%}\ \ WER};
        \node[font=\scriptsize, below=2pt of wer, text=AccAmber]
          {Word Error Rate — \textbf{Aceptable}};
      \end{tikzpicture}

      \vspace{8pt}
      \scriptsize
      \verde{Char. Accuracy: \textbf{89.37\,\%}} \quad
      \amber{Word Accuracy: \textbf{71.24\,\%}}
    \end{column}
  \end{columns}

  \vspace{6pt}
  \begin{alertblock}{Referencia industria}
    \scriptsize
    CER $<15\,\%$ = Bueno para documentos complejos con ruido.
    CER $<5\,\%$ = Producción. El pipeline alcanza \verde{\textbf{10.63\,\%}} sin GPU.
  \end{alertblock}
\end{frame}

\begin{frame}{Escala de Calidad OCR — Posicionamiento del Pipeline}
  \centering
  \begin{tikzpicture}
    % Barra de gradiente CER
    \shade[left color=AccGreen, right color=PriRed]
      (0,0) rectangle (10,0.7);
    % Marcas
    \foreach \x/\lbl in {0/0\%, 2/5\%, 4/10\%, 6/20\%, 8/30\%, 10/50\%+}
      \node[font=\tiny, below=2pt] at (\x,0) {\lbl};
    % Etiquetas sobre la barra
    \node[font=\tiny\bfseries, text=white] at (1.0,0.35) {Excelente};
    \node[font=\tiny\bfseries, text=white] at (3.2,0.35) {Muy bueno};
    \node[font=\tiny\bfseries, text=white] at (5.5,0.35) {Bueno};
    \node[font=\tiny\bfseries, text=white] at (7.5,0.35) {Regular};
    \node[font=\tiny\bfseries, text=white] at (9.2,0.35) {Deficiente};

    % Marcador: pipeline actual
    \draw[AccBlue, line width=2pt] (4.25,0) -- (4.25,1.5);
    \node[fill=AccBlue, text=white, font=\tiny\bfseries, rounded corners=3pt,
          inner sep=3pt, above] at (4.25,1.5) {Pipeline Actual\ 10.63\,\%};

    % Marcador: baseline estimado (sin procesamiento)
    \draw[PriRed, line width=1.5pt, dashed] (7.5,0) -- (7.5,1.5);
    \node[fill=PriRed, text=white, font=\tiny\bfseries, rounded corners=3pt,
          inner sep=3pt, above] at (7.5,1.5) {Baseline s/sello\ $\sim$30\,\%};

    % Marca objetivo
    \draw[AccGreen, dashed, line width=1pt] (2.0,-0.4) -- (2.0,0);
    \node[font=\tiny, text=AccGreen, below] at (2.0,-0.4) {Objetivo $<5\,\%$};
  \end{tikzpicture}

  \vspace{12pt}
  \begin{columns}[T]
    \begin{column}{0.31\textwidth}
      \scriptsize\centering
      \verde{\textbf{10.63\,\%\ CER}}\\Posición: Bueno\\Sin GPU requerida
    \end{column}
    \begin{column}{0.31\textwidth}
      \scriptsize\centering
      \amber{\textbf{28.76\,\%\ WER}}\\Impacto del sello\\en palabras completas
    \end{column}
    \begin{column}{0.31\textwidth}
      \scriptsize\centering
      \verde{\textbf{Mejora $>$5\,pp}}\\Objetivo cumplido\\vs.\ baseline estimado
    \end{column}
  \end{columns}
\end{frame}

\begin{frame}{Resultados Batch — 50 Imágenes Sintéticas (En Curso)}
  \begin{columns}[T]
    \begin{column}{0.50\textwidth}
      \scriptsize
      \begin{tabular}{@{}lrcc@{}}
        \toprule
        \textbf{Imagen} & \textbf{Tiempo} & \textbf{Conf.} & \textbf{Bloques} \\
        \midrule
        \rowcolor{AccGreen!10}
        syn\_0007 & 136\,s & 84.6\,\% & 47 \\
        \rowcolor{AccGreen!10}
        syn\_0024 & 138\,s & 36.7\,\% & 22 \\
        \rowcolor{AccGreen!10}
        syn\_0006 & 157\,s & 83.2\,\% & 125 \\
        syn\_0011 & 158\,s & 70.5\,\% & 76 \\
        syn\_0022 & 161\,s & 85.1\,\% & 106 \\
        syn\_0008 & 163\,s & 84.2\,\% & 217 \\
        syn\_0023 & 173\,s & 86.7\,\% & 231 \\
        syn\_0039 & 199\,s & 54.7\,\% & 18 \\
        syn\_0026 & 204\,s & 49.7\,\% & 43 \\
        syn\_0028 & 227\,s & 80.1\,\% & 137 \\
        \rowcolor{AccAmber!10}
        syn\_0035 & 233\,s & 90.2\,\% & 195 \\
        \rowcolor{PriRed!10}
        syn\_0001 & 239\,s & 78.1\,\% & 780 \\
        \midrule
        \textbf{Promedio} & \textbf{183\,s} & \textbf{73.7\,\%} & \textbf{164} \\
        \bottomrule
      \end{tabular}
      \vspace{4pt}
      \scriptsize\textit{$^\star$Tiempos incluyen primera descarga del modelo EasyOCR.\\Estimado real sin descarga: $\sim$40-50\,s/img.}
    \end{column}
    \begin{column}{0.47\textwidth}
      \scriptsize
      \begin{greenblock}{Estado del batch (12/50 completadas)}
        \begin{itemize}[leftmargin=1em, itemsep=3pt]
          \item \verde{12 imágenes OK} — 0 fallos
          \item RAM pico: \textbf{85.9\,\%} (2.2\,GB libre)
          \item Workers activos: \textbf{8 paralelos}
          \item Guardado incremental activo
        \end{itemize}
      \end{greenblock}
      \vspace{6pt}
      \begin{blueblock}{Distribución de confianza (12 imgs)}
        \begin{itemize}[leftmargin=1em, itemsep=3pt]
          \item Alta $\geq 80\,\%$:\ \ \verde{\textbf{7 imágenes (58\,\%)}}
          \item Media $50-80\,\%$:\ \amber{\textbf{3 imágenes (25\,\%)}}
          \item Baja $< 50\,\%$:\ \ \rojo{\textbf{2 imágenes (17\,\%)}}
        \end{itemize}
      \end{blueblock}
      \vspace{6pt}
      \scriptsize\rojo{\textbf{Nota:}} syn\_0024 (36.7\,\%) y syn\_0026 (49.7\,\%)
      son imágenes con sello muy agresivo o texto muy pequeño.
    \end{column}
  \end{columns}
\end{frame}

\begin{frame}{Comparativa Teórica: Propuesta vs.\ Implementación}
  \centering
  \footnotesize
  \begin{tabular}{@{}p{3.2cm}cc@{}}
    \toprule
    \textbf{Métrica}
      & \original{Propuesta Original}
      & \mejora{Implementación Actual} \\
    \midrule
    \rowcolor{RowA}
    CER objetivo
      & $<$ 5\,\% (ambicioso)
      & \verde{10.63\,\%} (\textit{bueno}) \\[3pt]
    \rowcolor{RowB}
    WER obtenido
      & No definido explícitamente
      & \amber{28.76\,\%} (\textit{aceptable}) \\[3pt]
    \rowcolor{RowA}
    GPU requerida
      & \rojo{Sí} (LaMa $>4$\,GB VRAM)
      & \verde{No} (100\,\% CPU) \\[3pt]
    \rowcolor{RowB}
    Anotaciones manuales
      & \rojo{Sí} (YOLO bboxes)
      & \verde{No} (sin supervisión) \\[3pt]
    \rowcolor{RowA}
    Tiempo por imagen
      & $\sim$5\,s con GPU
      & $\sim$40-50\,s CPU (sin descarga) \\[3pt]
    \rowcolor{RowB}
    Sellos translúcidos
      & \rojo{Difícil} (YOLO + HSV fijo)
      & \verde{Sí} (clasificación cromática) \\[3pt]
    \rowcolor{RowA}
    Portabilidad Colab
      & \rojo{Limitada} (GPU tier)
      & \verde{Completa} (CPU gratuito) \\[3pt]
    \rowcolor{RowB}
    Imágenes procesadas
      & 100 (benchmark humano)
      & \verde{50+ sintéticas + reales} \\[3pt]
    \rowcolor{RowA}
    Post-corrección NLP
      & pyspellchecker básico
      & \verde{Levenshtein + NER spaCy} \\
    \bottomrule
  \end{tabular}
\end{frame}

% ══════════════════════════════════════════════════════════════════════════════
\section{Análisis Cuantitativo}
% ══════════════════════════════════════════════════════════════════════════════

\begin{frame}{Distribución de Confianza OCR — Batch 50 imágenes}
  \begin{columns}[T]
    \begin{column}{0.52\textwidth}
      \centering\vspace{4pt}
      \begin{tikzpicture}
        \begin{axis}[
          width=6.5cm, height=5.5cm,
          ybar,
          bar width=14pt,
          xlabel={\scriptsize Rango de confianza},
          ylabel={\scriptsize Imágenes},
          xtick={1,2,3,4,5},
          xticklabels={$<$40\%, 40-60\%, 60-80\%, 80-90\%, $>$90\%},
          xticklabel style={font=\tiny, rotate=30, anchor=east},
          ymin=0, ymax=6,
          ytick={0,1,2,3,4,5,6},
          yticklabel style={font=\tiny},
          axis lines=left,
          ymajorgrids=true,
          grid style={dashed, gray!30},
          nodes near coords,
          nodes near coords style={font=\tiny},
        ]
          \addplot[fill=PriRed!70, draw=PriRed] coordinates {(1,2)};
          \addplot[fill=AccAmber!80, draw=AccAmber] coordinates {(2,2)};
          \addplot[fill=AccAmber!50, draw=AccAmber!80] coordinates {(3,1)};
          \addplot[fill=AccGreen!70, draw=AccGreen] coordinates {(4,5)};
          \addplot[fill=AccGreen!90, draw=AccGreen!120] coordinates {(5,2)};
        \end{axis}
      \end{tikzpicture}
      \vspace{-4pt}
      \scriptsize\textit{Basado en 12 imágenes completadas del batch}
    \end{column}
    \begin{column}{0.45\textwidth}
      \vspace{8pt}
      \begin{itemize}[leftmargin=1em, itemsep=8pt]
        \item \verde{\textbf{58\,\%}} de imágenes con confianza $\geq 80\,\%$ —
              OCR de alta calidad
        \item \amber{\textbf{25\,\%}} en rango medio (50-80\,\%) —
              texto con ruido o sello parcial
        \item \rojo{\textbf{17\,\%}} confianza baja — sellos muy
              agresivos o texto muy pequeño
        \item Promedio global: \amber{\textbf{73.7\,\%}}
        \item Máximo alcanzado: \verde{\textbf{90.2\,\%}} (syn\_0035)
        \item Mínimo: \rojo{\textbf{36.7\,\%}} (syn\_0024)
      \end{itemize}
    \end{column}
  \end{columns}
\end{frame}

\begin{frame}{Análisis de Tiempos por Fase (Estimados CPU)}
  \begin{columns}[T]
    \begin{column}{0.50\textwidth}
      \centering\vspace{4pt}
      \begin{tikzpicture}
        \begin{axis}[
          xbar, bar width=12pt,
          width=6.5cm, height=5cm,
          xlabel={\scriptsize Segundos (promedio)},
          ytick={1,2,3},
          yticklabels={Consolidar, Fase 4 OCR, Fase 3A Sello},
          yticklabel style={font=\scriptsize},
          xmin=0, xmax=200,
          nodes near coords,
          nodes near coords style={font=\tiny},
          axis lines=left,
          xmajorgrids=true,
          grid style={dashed,gray!30},
        ]
          \addplot[fill=AccTeal!70, draw=AccTeal] coordinates {(10,1)};
          \addplot[fill=AccAmber!80, draw=AccAmber] coordinates {(165,2)};
          \addplot[fill=OldRed!70, draw=OldRed] coordinates {(8,3)};
        \end{axis}
      \end{tikzpicture}
      \scriptsize\textit{$^\star$Tiempos del primer batch incluyen\\descarga EasyOCR ($\sim$140\,s extra)}
    \end{column}
    \begin{column}{0.47\textwidth}
      \vspace{8pt}
      \begin{blueblock}{Breakdown real (primer batch)}
        \scriptsize
        \begin{tabular}{lr}
          Fase 3A (inpainting) & $\sim$8\,s \\
          EasyOCR init + inferencia & $\sim$155\,s \\
          Spell-check + NER & $\sim$8\,s \\
          Consolidación + I/O & $\sim$10\,s \\
          \midrule
          \textbf{Total (con modelo)} & $\sim$183\,s \\
          \textbf{Total (modelo cacheado)} & $\sim$40-50\,s \\
        \end{tabular}
      \end{blueblock}
      \vspace{6pt}
      \begin{exampleblock}{Bottleneck identificado}
        \scriptsize
        \textbf{EasyOCR} consume $>$85\,\% del tiempo total.
        Mejora obvia: GPU (reducción $\sim$10$\times$) o batching de imágenes.
      \end{exampleblock}
    \end{column}
  \end{columns}
\end{frame}

% ══════════════════════════════════════════════════════════════════════════════
\section{Conclusiones}
% ══════════════════════════════════════════════════════════════════════════════

\begin{frame}{Logros vs.\ Objetivos Originales}
  \centering
  \footnotesize
  \begin{tabular}{@{}p{4.0cm}ccc@{}}
    \toprule
    \textbf{Objetivo}
      & \textbf{Meta}
      & \textbf{Resultado}
      & \textbf{Estado} \\
    \midrule
    Reducción CER $\geq 5$\,pp
      & $>5$\,pp
      & $\sim$20\,pp estimados
      & \verde{\checkmark} \\[3pt]
    CER $< 15$\,\% (bueno)
      & $<15$\,\%
      & \textbf{10.63\,\%}
      & \verde{\checkmark} \\[3pt]
    100\,\% CPU (sin GPU)
      & Portabilidad
      & Colab gratuito
      & \verde{\checkmark} \\[3pt]
    Sin anotaciones manuales
      & Zero-label
      & Clasificación cromática
      & \verde{\checkmark} \\[3pt]
    Manejo sellos translúcidos
      & Caso difícil
      & 5 categorías RGB
      & \verde{\checkmark} \\[3pt]
    Dataset $\geq 100$ docs
      & 100 docs
      & 200 sintéticas + reales
      & \verde{\checkmark} \\[3pt]
    CER $<5$\,\% (producción)
      & $<5$\,\%
      & 10.63\,\% (gap $\sim$6\,pp)
      & \amber{$\sim$} \\
    \bottomrule
  \end{tabular}

  \vspace{10pt}
  \begin{columns}[T]
    \begin{column}{0.32\textwidth}
      \centering\verde{\LARGE\checkmark}\\
      \scriptsize 6 de 7 objetivos\\completamente logrados
    \end{column}
    \begin{column}{0.32\textwidth}
      \centering\amber{\Large$\sim$}\\
      \scriptsize CER $<5\,\%$ requiere\\GPU o más datos GT
    \end{column}
    \begin{column}{0.32\textwidth}
      \centering\azul{\large$\uparrow$}\\
      \scriptsize Pipeline listo para\\escalar a producción
    \end{column}
  \end{columns}
\end{frame}

\begin{frame}{Trabajo Futuro y Próximos Pasos}
  \begin{columns}[T]
    \begin{column}{0.48\textwidth}
      \begin{blueblock}{Mejoras al Pipeline}
        \begin{itemize}[leftmargin=1em, itemsep=5pt]
          \item \textbf{GPU acceleration:} EasyOCR con CUDA —
                reducción $\sim$10$\times$ en tiempo
          \item \textbf{LaMa opcional:} Para imágenes con sello
                muy agresivo (GPU tier)
          \item \textbf{Fine-tuning OCR:} EasyOCR en documentos
                legales españoles
          \item \textbf{Umbral adaptativo:} Ajuste automático
                de parámetros cromáticos por tipo de documento
        \end{itemize}
      \end{blueblock}
    \end{column}
    \begin{column}{0.48\textwidth}
      \begin{greenblock}{Expansión del Benchmark}
        \begin{itemize}[leftmargin=1em, itemsep=5pt]
          \item Completar batch de \textbf{50 imgs sintéticas}
                (en curso — 12/50 procesadas)
          \item Generar ground truth para \textbf{20 docs reales}
          \item Medir CER \textbf{antes y después} para calcular
                mejora exacta del pipeline
          \item Comparar con \textbf{Tesseract baseline} como
                punto de referencia externo
          \item Publicar dataset sintético como contribución abierta
        \end{itemize}
      \end{greenblock}
    \end{column}
  \end{columns}
\end{frame}

\begin{frame}{Resumen de Contribuciones}
  \begin{columns}[T]
    \begin{column}{0.55\textwidth}
      \begin{enumerate}[label=\rojo{\textbf{\arabic*.}}, leftmargin=1.5em, itemsep=8pt]
        \item \textbf{Clasificación cromática sin supervisión}\\
              5 categorías de píxel por propiedades físicas del color,
              sin necesidad de anotaciones ni GPU.

        \item \textbf{Horizontal/Block Sweep Filter}\\
              Coherencia contextual para evitar falsos positivos
              en regiones de texto denso.

        \item \textbf{Inpainting híbrido adaptativo}\\
              Selección automática TELEA/Navier-Stokes según
              densidad de la máscara.

        \item \textbf{Post-corrección NLP con protección NER}\\
              Levenshtein $\leq 2$ sin corromper entidades nombradas
              (personas, instituciones, fechas).

        \item \textbf{Pipeline 100\,\% CPU, Colab-ready}\\
              Ejecutable en Google Colab gratuito sin configuración local.
      \end{enumerate}
    \end{column}
    \begin{column}{0.42\textwidth}
      \centering\vspace{10pt}
      \begin{tikzpicture}
        \node[draw=AccGreen, line width=2pt, rounded corners=8pt,
              fill=AccGreen!10, minimum width=5cm, minimum height=2.5cm,
              align=center] {
          \Large\bfseries\verde{10.63\,\%}\\[4pt]
          \normalsize CER alcanzado\\[6pt]
          \scriptsize sin GPU $\cdot$ sin anotaciones\\
          \scriptsize sin configuración local
        };
      \end{tikzpicture}

      \vspace{12pt}
      \begin{tikzpicture}
        \node[draw=AccAmber, line width=1.5pt, rounded corners=8pt,
              fill=AccAmber!10, minimum width=5cm, minimum height=1.6cm,
              align=center] {
          \large\bfseries\amber{89.4\,\%}\\[3pt]
          \scriptsize Precisión a nivel de carácter
        };
      \end{tikzpicture}
    \end{column}
  \end{columns}
\end{frame}

% ─── CIERRE ───────────────────────────────────────────────────────────────────
\begin{frame}[standout]
  \centering
  {\Large\bfseries Gracias}\\[16pt]
  {\normalsize Detección y Eliminación Inteligente de Sellos\\
   para Mejora del OCR en Documentos Oficiales}\\[20pt]
  \begin{tabular}{rl}
    CER alcanzado: & \textbf{10.63\,\%} \\
    Precisión char: & \textbf{89.37\,\%} \\
    GPU requerida: & \textbf{No} \\
    Anotaciones: & \textbf{Ninguna} \\
  \end{tabular}\\[20pt]
  {\small Manuel Cruz Garrote}\\
  {\small Universidad del Rosario $\cdot$ Mayo 2026}
\end{frame}

% ─── REFERENCIAS ──────────────────────────────────────────────────────────────
\begin{frame}[allowframebreaks]{Referencias}
  \scriptsize
  \begin{enumerate}[label={[\arabic*]}, leftmargin=1.8em, itemsep=4pt]
    \item Suvorov, R. et al.\ (2022). Resolution-robust Large Mask Inpainting
          with Fourier Convolutions. \textit{WACV 2022}.
    \item Du, Y. et al.\ (2023). PP-OCRv4: Self-Supervised Pre-Training for
          Text Detection and Recognition. \textit{arXiv:2309.06383}.
    \item Jocher, G. et al.\ (2023). Ultralytics YOLOv8. \textit{GitHub}.
    \item Baek, J. et al.\ (2019). CRAFT: Character-Region Awareness for
          Text Detection. \textit{CVPR 2019}.
    \item Liang, J. et al.\ (2022). An adaptive deskewing algorithm for
          document images. \textit{Pattern Recognition Letters}.
    \item Terven, J. \& Cordova, D. (2023). A Comprehensive Review of YOLO
          Architectures. \textit{arXiv:2304.00501}.
    \item Smith, R. (2021+). Tesseract OCR Engine v5. \textit{GitHub}.
    \item Honnibal, M. \& Montani, I. (2023). spaCy: Industrial-strength NLP.
          \textit{explosion.ai}.
  \end{enumerate}
\end{frame}

\end{document}
